% Modernised reading — fonds Alexandre Grothendieck, shelfmark 7, pages 1-71.
%
% This is an INTERPRETATION, not a transcription. It re-reads the pages in
% today's notation and vocabulary, states the hypotheses the manuscript leaves
% implicit, and drops the critical apparatus so that the mathematics reads
% straight through. Everything the brackets and underlines carried — a doubtful
% reading, a notation he used differently, a missing passage — moves into a
% footnote. It is held to being mathematically correct as it stands; in any
% dispute of substance the transcription governs: whoever cites Grothendieck
% must cite the batch-NN.fr.tex files.
%
% Scope: the folder was read WHOLE — all four batches, pages 1-71 — before any
% of this was written, and it is written as ONE file for the shelfmark. Pages
% 1 and 19 are blank. Two dated layers: pp. 20-71 are the talk notes
% (« laïus ») of November-December 1966 (dates 15.11.66 p. 35, 6.12.1966
% p. 56, « Nov. Déc. 1966 » p. 20, all checked on the facsimile by Michel Hua);
% pp. 2-18 are dated « Juillet-Août 1967 », the 1967 read by elimination. The
% reading follows page order, so the later layer comes first; the spine
% section says so.
%
% Notation fixed for the whole folder (see « Conventions »):
%   X^{nu+1} the (nu+1)-fold fibre product over S (his X^{[nu]} on p. 2 has
%   nu+1 factors: his own margin on p. 9 warns of the shift); Delta^{(nu)}(i)
%   the i-th infinitesimal neighbourhood of the diagonal; X^{nu+1}/X the
%   formal completion; P^nu its ring (principal parts), Q^nu = P^{nu+1};
%   (X/S)_inf for what the folder calls the « site cristallin » — NO divided
%   powers appear anywhere in the folder, so today's name is the
%   infinitesimal site; (X/S)_strat for the stratifying site; C^*(V) the
%   Čech-Alexander complex of a stratified module.
%
% Substantive departures from the pages, each footnoted where it occurs:
%   - p. 13-15: the formal Poincaré lemma needs S of characteristic 0 (the
%     word is illegible on p. 13; p. 15 says it); it is false in char. p;
%   - p. 21: the comparison theorem needs X smooth (the parenthesis says only
%     « loc. de type fini »);
%   - p. 40: the overconvergence bound is |a_n| <= C rho^{-n}, not O(1) rho^n;
%   - p. 41: u : A -> A' induces H(Omega(A)) -> H(Omega(A')), and the
%     independence of the lift is stated after tensoring with K (the integral
%     form is his hope, and p. 42 says W-M prove only the rational one);
%   - p. 58: the boxed formula is Q^{i-1}[M] = P^i (x) M (index);
%   - p. 70: the argument for H_inf -> H_strat iso is elliptic; restated
%     with what it actually needs.
% Announced and never established in the folder: the conjecture of pp. 55
% and 69 (stratifying cohomology of a singular complex variety); the « Pb »
% of pp. 70-71; §7 and §8 of the plan of p. 59 (unless pp. 2-18 are them, a
% hypothesis of ours); the « appendice » of p. 51; tests a) b) of p. 18.
%
% Pass: Opus 5.5 (claude-opus-5-5), 2026-10-03 — first pass, unchecked against
% the pages by a human.
\documentclass[11pt,a4paper]{article}
\input{../preamble/grothendieck}

\folder{7}
\pages{1}{71}
\dating{1966-1967}
\watermark{Édition de démonstration\\interprétation personnelle de l'œuvre}
\foldertitle{Notes laïus novembre-décembre 1967 (Cristaux) : notes manuscrites (1966-1967) — lecture modernisée du dossier entier}

\begin{document}

\section*{Résumé}

\begin{resume}
Ces feuillets sont les notes d'un mathématicien qui prépare une série
d'exposés — il les appelle son « laïus » — sur une question précise : comment
compter les « trous » d'une variété définie par des équations, quand ces
équations sont écrites non pas sur les nombres réels ou complexes, mais sur un
corps où $p \cdot 1 = 0$ pour un nombre premier $p$.

Sur les nombres complexes, la question a une réponse ancienne et belle. Une
forme différentielle fermée qui n'est pas la différentielle d'une fonction
signale un trou : sur le plan privé de l'origine, la forme « angle » $d\theta$
est fermée, mais $\theta$ n'est pas une fonction bien définie, parce qu'on peut
tourner autour du trou. Compter ces formes, à l'exacte près, c'est la
\emph{cohomologie de De Rham}. Le dossier part de ce fait, et de sa version
algébrique : les formes différentielles s'écrivent avec des polynômes, et la
cohomologie qu'on en tire donne les bons nombres de trous (pages 20 à 22).

Le problème est qu'en caractéristique $p$ cette recette se dérègle. La
dérivée de $t^{p}$ est $p\,t^{p-1} = 0$ : il y a beaucoup trop de fonctions
« constantes », et les nombres qu'on obtient sont faux. On possède bien
d'autres théories, dites $\ell$-adiques, qui marchent pour tout nombre premier
$\ell$ \emph{différent} de $p$ ; mais il manque celle qui verrait $p$
lui-même, et le dossier explique pourquoi on ne peut s'en passer (pages 26 à
32). Une première tentative, celle de Washnitzer et Monsky, consiste à
« relever » la variété en caractéristique nulle et à y prendre la cohomologie
de De Rham ; le dossier en fait la critique (pages 34 à 46).

L'idée qui arrive est la suivante. Au lieu de travailler avec des formes
différentielles, qui exigent de savoir dériver, on regarde comment un objet se
\emph{prolonge} aux voisinages infiniment petits de la variété : les points
« infiniment proches » les uns des autres, rangés par ordre de proximité. Un
objet qui se prolonge à tous ces voisinages, de façon cohérente et rigide,
s'appelle aujourd'hui un \emph{cristal} — un cristal croît en gardant
exactement sa structure. La cohomologie d'un site formé de ces voisinages
redonne la cohomologie de De Rham en caractéristique nulle, sans jamais
écrire une forme différentielle (pages 53 à 55) ; et elle se calcule par un
procédé à la Čech, avec les fonctions sur les voisinages formels des
diagonales, ce qui occupe toute la fin du dossier (pages 56 à 71). La couche
la plus tardive, datée de l'été 1967 (pages 2 à 18), refait la construction
proprement et retourne à la caractéristique $p$ — pour constater que, telle
quelle, elle y échoue.

Ce qui rend ces pages précieuses, au-delà de leur mathématique, c'est qu'on y
voit un espoir se former et se défaire. « Mon espoir initial était qu'il soit
“le” bon candidat », écrit-il à la page 16, et deux lignes plus bas : « C'est
faux malheureusement. » Une remarque reçoit après coup, en marge, un « C'est
sûrement faux ! » ; une relation entre deux idées est dite « heuristique et
sentimentale », faute d'être comprise (page 52) ; un exposé s'ouvre sur le
projet de « faire sentir un complexe de problèmes » plutôt que d'énoncer des
théorèmes (page 20). On y apprend qu'un programme de recherche, chez lui,
commence par un inventaire honnête de ce qui ne marche pas.

Où cela mène : la réponse en caractéristique $p$ viendra des \emph{puissances
divisées}, absentes de ces pages — c'est la cohomologie cristalline, construite
par Berthelot ; le site plus gros que propose la page 17 a des cousins dans le
site convergent d'Ogus et la cohomologie rigide de Berthelot. Les noms à
chercher sont ci-dessous.

\keywords{algebraic de Rham cohomology, comparison theorem, Hodge-to-de Rham spectral sequence, Leray spectral sequence, l-adic cohomology, independence of l, p-adic cohomology, Witt vectors, flat cohomology, Monsky-Washnitzer cohomology, overconvergent power series, weakly complete algebra, Gauss-Manin connection, Cartan homotopy formula, stratification, crystal, infinitesimal site, stratifying site, Čech-Alexander complex, Alexander-Spanier cochains, principal parts, linearization of differential operators, complexes of differential operators, formal Poincaré lemma, topological invariance, GAGA, crystalline cohomology, convergent site}
\end{resume}

\pagerange{2}{3}
\section*{Le fil du dossier, et les conventions}

\subsection*{Deux couches, et l'ordre des feuillets}

Le dossier contient deux couches de dates différentes, rangées par
l'inventaire dans l'ordre inverse de leur rédaction.\footnote{Les pages 1 et
19 sont blanches. La date de la page 2, « Juillet-Août 1967 », a un troisième
chiffre à peine formé : la transcription le lit 1967 par élimination (6 ou 8
sont possibles, 1987 est exclu par le contenu). Les dates de la couche de
1966 ont été vérifiées sur le fac-similé : « Nov.\ Déc.\ 1966 » à la page 20,
le dernier chiffre corrigé en rouge ; « 15.11.66 » à la page 35, d'abord écrit
« 25 » ; « 6.12.1966 » à la page 56. Le titre du catalogue, « novembre-décembre
1967 », ne s'accorde donc pas avec les dates portées sur les feuillets de
l'exposé ; nous suivons les feuillets.}

\begin{itemize}
\item \emph{Pages 20 à 71 : les notes du « laïus », novembre--décembre 1966.}
  Un premier exposé paginé par lui de 1 à 15 (pages 20 à 34) : §\,1, la
  cohomologie de De Rham ; §\,2, le tour d'horizon $\ell$-adique ; §\,3, les
  essais de cohomologie $p$-adique, jusqu'au début de Washnitzer--Monsky. Une
  reprise datée du 15 novembre (pages 35 à 52) : Washnitzer--Monsky en
  détail, puis la relation heuristique avec Gauss--Manin. Un §\,4 sur les
  topos stratifiant et cristallin (pages 53 à 55). Une reprise datée du
  6 décembre (pages 56 à 58), puis un plan d'ensemble en huit paragraphes
  (page 59) et le §\,6 de ce plan, « Calculs Čechistes » (pages 59 à 71).
\item \emph{Pages 2 à 18 : « Anneaux simpliciaux canoniques et modules
  stratifiés », été 1967.} Six paragraphes numérotés qui reprennent, de façon
  systématique, la construction esquissée aux pages 56 à 58, et la poussent
  jusqu'au lemme de Poincaré formel (pages 2 à 15) ; puis des « Commentaires
  sur la définition d'une cohomologie $p$-adique en car.\ $p > 0$ »
  (pages 16 à 18).
\end{itemize}

Le plan de la page 59 annonce un §\,7, « Comparaison de la cohomologie de De
Rham et cristalline », et un §\,8, « Retour sur le cas de la car.\ $p$ », que
la couche de 1966 n'écrit pas. Les pages 2 à 15 et 16 à 18 en ont exactement
la matière ; c'est une hypothèse de notre part que ce soient eux, et rien sur
les feuillets ne le dit.\footnote{La concordance des dates et des sujets avec
les exposés de Grothendieck à l'IHÉS de novembre--décembre 1966, rédigés par
J.~Coates et O.~Jussila sous le titre « Crystals and the De Rham cohomology
of schemes » et publiés dans \emph{Dix exposés sur la cohomologie des
schémas} (1968), suggère que ces pages en sont les notes préparatoires. C'est
une identification de notre part, que les feuillets ne font pas.}

\subsection*{Conventions, valables pour tout le dossier}

On fixe une base $S$ et un préschéma relatif $f : X \to S$ (aujourd'hui : un
schéma au-dessus de $S$).

\begin{itemize}
\item $X^{\nu+1} = X \times_{S} \cdots \times_{S} X$ ($\nu+1$ facteurs), et
  $\Delta^{(\nu)}(i) \subset X^{\nu+1}$ le $i$-ème voisinage infinitésimal de
  la diagonale. On note $X^{\nu+1}/X$ le complété formel de $X^{\nu+1}$ le
  long de la diagonale, c'est-à-dire la limite inductive des
  $\Delta^{(\nu)}(i)$ ; tous ont pour espace sous-jacent celui de
  $X$.\footnote{La page 2 note $X^{[n]}$ le produit à $n+1$ facteurs, et son
  complété $X^{\widehat{[n]}}$ ; les pages 59 et 60 écrivent $X^{\nu+1}$ et
  $X^{(\nu+1)}/X$. Le décalage d'une unité entre l'exposant et le nombre de
  facteurs est la source d'erreur qu'il se signale lui-même en marge de la
  page 9 (« Mention $+1$ ! »). On compte ici les facteurs, partout.}
\item $\mathcal{P}^{\nu} = \varprojlim_{i} \mathcal{O}_{\Delta^{(\nu)}(i)}$,
  vu comme pro-objet de la catégorie des faisceaux d'anneaux sur l'espace de
  $X$ : c'est le faisceau des \emph{parties principales} d'ordre infini.
  Pour $\nu$ variable, les $\mathcal{P}^{\nu}$ forment un anneau cosimplicial
  $\mathcal{P}^{*}$, avec $\mathcal{P}^{0} = \mathcal{O}_X$. On pose
  $\mathcal{Q}^{\nu} = \mathcal{P}^{\nu+1}$ (le décalé).
\item Un $\mathcal{O}_X$-Module $M$ donne $\mathcal{P}^{\nu}[M] =
  \mathcal{P}^{\nu} \otimes_{\mathcal{O}_X} M$ et $\mathcal{Q}^{\nu}[M] =
  \mathcal{Q}^{\nu} \otimes_{\mathcal{O}_X} M$, le produit tensoriel étant
  toujours pris pour la structure de $\mathcal{O}_X$-algèbre venue de la
  \emph{dernière} projection (« extrême-droite »). Quand on passe à la limite
  projective d'un tel pro-objet, on écrit $\varprojlim$.
\item $\mathrm{Mod}(X/S)$ désigne la catégorie dont les objets sont les
  $\mathcal{O}_X$-Modules et les morphismes les \emph{opérateurs
  différentiels} relatifs à $S$ ; c'est le sens que la page 58 lui donne en
  toutes lettres. Elle est additive, non abélienne.
\item Le dossier dit « site cristallin » pour le site des épaississements
  nilpotents \emph{sans puissances divisées}. Rien, dans ces pages,
  n'introduit de puissances divisées ; on écrit donc $(X/S)_{\mathrm{inf}}$,
  le \emph{site infinitésimal}, là où il écrit
  $X_{\mathrm{cris}}$.\footnote{Dans la terminologie fixée ensuite —
  « Crystals and the De Rham cohomology of schemes » (1968), puis la thèse de
  Berthelot, \emph{Cohomologie cristalline des schémas de caractéristique
  $p > 0$} (1974) — le site sans puissances divisées s'appelle site
  infinitésimal, et « cristallin » est réservé au site à puissances
  divisées. Le nom ne change rien en caractéristique nulle, où les deux sites
  ont même cohomologie ; il change tout en caractéristique $p$, et c'est
  justement l'échec constaté à la page 16.} Le \emph{site stratifiant}
  $(X/S)_{\mathrm{strat}}$ est le sous-site formé des épaississements
  $U \hookrightarrow T$ pour lesquels l'inclusion de $U$ dans $X$ se
  prolonge, au moins localement, en un $S$-morphisme $T \to
  X$.\footnote{Le dossier ne donne pas la définition par écrit : la page 54
  l'annonce (« Définition du site et du topos stratifiant ») et la page 59
  invoque « la restriction sur ces objets du site stratifiant parmi ceux du
  site cristallin ». Nous la restituons sous la forme qui rend vraies les
  assertions des pages 53, 54 et 59 : avec elle, la donnée d'un Module sur
  ce site, à flèches de transition des isomorphismes, est une stratification,
  et le faisceau $\widetilde{X}$ représenté par $X$ lui-même couvre l'objet
final.}
\item Un Module stratifié se note $\underline{V}$. Son \emph{complexe de
  Čech--Alexander} est le faisceau cosimplicial sur $X_{\mathrm{zar}}$
  \[
    C^{\nu}(\underline{V}) = \varprojlim_{i} \underline{V}_{\Delta^{(\nu)}(i)} ,
  \]
  l'image inverse par n'importe laquelle des $\nu+1$ projections, toutes
  canoniquement isomorphes grâce à la stratification ; la page 12 l'appelle
  $C^{\cdot}_{\mathrm{str}}(\underline{V})$, la page 56
  $C^{*}(\mathcal{F})$, les pages 61 à 70 $\mathcal{F}^{*}$.
\item $\mathbb{H}$ désigne l'hypercohomologie. $k$ est un corps parfait de
  caractéristique $p > 0$, $W = W(k)$ l'anneau de ses vecteurs de Witt,
  $K = W[1/p]$ son corps des fractions, $W_n = W/p^{n}W$.
\end{itemize}

\pagerange{2}{3}
\section*{Anneaux cosimpliciaux canoniques (pages 2 et 3)}

La rédaction de l'été 1967 s'ouvre sur un avertissement : formellement, ce
qui suit ne dépend pas des « histoires de sites cristallins ou
stratifiants ».\footnote{« stratifiants » est une lecture incertaine ; le
sens n'en dépend pas.} C'est de l'algèbre cosimpliciale sur l'espace de $X$,
et c'est ce qui permet ensuite d'en tirer des sites.

Pour $\nu$ variable, les complétés $X^{\nu+1}/X$ forment un schéma formel
simplicial — le nerf de Čech de $X \to S$, complété le long de la diagonale
— dont tous les termes ont l'espace de $X$. Leurs faisceaux structuraux
forment donc l'anneau cosimplicial $\mathcal{P}^{*}$ des conventions, et
$\mathcal{P}^{0} = \mathcal{O}_X$. Ce n'est pas une $\mathcal{O}_X$-algèbre
cosimpliciale : les deux cofaces $\mathcal{O}_X = \mathcal{P}^{0}
\rightrightarrows \mathcal{P}^{1}$ sont en général distinctes, et mettent sur
$\mathcal{P}^{1}$ deux structures de $\mathcal{O}_X$-algèbre
\emph{différentes} — la gauche et la droite. Tout ce qui suit vit de cet
écart.

Le \emph{décalé} $\mathcal{Q}^{*} = \mathcal{S}\mathcal{P}^{*}$ s'obtient en
ajoutant un point : si l'on voit $\mathcal{P}$ comme foncteur sur les
ensembles finis non vides, $\mathcal{Q}(I) = \mathcal{P}(I \sqcup e)$, avec
$e$ un ensemble à un élément, d'où $\mathcal{Q}^{\nu} = \mathcal{P}^{\nu+1}$.
L'inclusion $I \to I \sqcup e$ donne un morphisme d'anneaux cosimpliciaux
\[
  \alpha : \mathcal{P}^{*} \longrightarrow \mathcal{Q}^{*} ,
  \qquad \alpha^{\nu} : \mathcal{P}^{\nu} \to \mathcal{P}^{\nu+1} .
\]
Exemple : $S = \operatorname{Spec} k$, $X = \operatorname{Spec} A$, la
diagonale $X \to X \times_S X$ étant une immersion nilpotente (par exemple
$X$ radiciel sur $S$), de sorte qu'aucune complétion n'est
nécessaire.\footnote{L'hypothèse est en partie illisible à la page 3
(« immersion nil-[\ldots] ») ; « nilpotente » est la seule lecture qui rende
juste la formule qui suit.} Alors
\[
  \Gamma \mathcal{P}^{\nu} = A^{\otimes_k (\nu+1)} ,
  \qquad
  \alpha^{\nu}(x_0 \otimes \cdots \otimes x_\nu)
  = x_0 \otimes \cdots \otimes x_\nu \otimes 1 ,
\]
les opérations cosimpliciales étant celles du nerf de Čech de $k \to A$.

\pagerange{4}{8}
\section*{Modules, et la linéarisation des opérateurs différentiels (pages 4 à 8)}

Soit $M$ un $\mathcal{O}_X$-Module. On pose $\mathcal{P}^{\nu}[M] =
\mathcal{P}^{\nu} \otimes_{\mathcal{O}_X} M$, pour la structure
extrême-droite. Par définition c'est un pro-objet ; on peut le remplacer par
sa limite, un vrai faisceau, sous des conditions de finitude — $X$ localement
noethérien, $\Omega^{1}_{X/S}$ et $M$ de type fini — qui rendent ce passage à
la limite inoffensif.\footnote{La page 4 est très lacunaire à cet endroit ;
les conditions sont celles qu'on peut lire entre les mots illisibles, et
leur rôle (remplacer les profaisceaux par leur limite) est explicite. Nous ne
prétendons pas les avoir restituées entièrement.} C'est un
$\mathcal{P}^{\nu}$-Module, fonctoriel en $M$.

Le point est dans les $\mathcal{Q}^{\nu}[M] = \mathcal{Q}^{\nu}
\otimes_{\mathcal{O}_X} M$. Les opérations cosimpliciales de $\mathcal{Q}^{*}$
ne touchent pas au dernier facteur, celui du point ajouté ; elles sont donc
linéaires pour la structure extrême-droite, et les $\mathcal{Q}^{\nu}[M]$
forment un \emph{Module cosimplicial} $\mathcal{Q}^{*}[M]$ sur l'anneau
cosimplicial $\mathcal{Q}^{*}$.

\emph{Fonctorialité en opérateurs différentiels.} Le terme
$\mathcal{Q}^{0}[M] = \mathcal{P}^{1} \otimes_{\mathcal{O}_X} M$ est le
pro-Module des parties principales de $M$. Un opérateur différentiel
$D : M \to N$ d'ordre $\leq n$ est, par définition, la même chose qu'une
application $\mathcal{O}_X$-linéaire $\mathcal{P}^{1}_{(n)} \otimes M \to N$ ;
il induit donc un morphisme
\[
  \mathcal{Q}^{0}[D] : \mathcal{Q}^{0}[M] \longrightarrow \mathcal{Q}^{0}[N]
\]
qui est linéaire pour la structure \emph{gauche}, et respecte la
stratification, mais non la structure extrême-droite. Ainsi $M \mapsto
\mathcal{Q}^{0}[M]$ est un foncteur de $\mathrm{Mod}(X/S)$ vers les
pro-Modules stratifiés, et $M$ s'en récupère comme sous-faisceau des sections
horizontales, $D$ comme la restriction de $\mathcal{Q}^{0}[D]$ à ces
sections.\footnote{C'est ce qu'on appelle aujourd'hui la
\emph{linéarisation} des opérateurs différentiels : EGA~IV, 16.8, pour la
correspondance entre opérateurs différentiels et applications linéaires sur
les parties principales, et Berthelot (1974), où le foncteur $L$ joue le même
rôle sur le site cristallin. Les pages 6 et 7 sont très lacunaires : la
récupération de $M$ par les sections horizontales est lisible, l'énoncé
exact sur $D$ l'est à peine. La mention « (Artin-Rees) » de la page 6, lue
sous réserve, porte sans doute sur le contrôle du passage à la limite dans
les pro-Modules.}

\emph{Une résolution canonique.} On obtient autrement $\mathcal{Q}^{*}$ en
regardant
\[
  X^{3}/X \rightrightarrows X^{2}/X \xrightarrow{\ \mathrm{pr}_2\ } X ,
\]
c'est-à-dire en formant le nerf de Čech, au-dessus de $X$, du morphisme
$\mathrm{pr}_2 : X^{2}/X \to X$. Ce morphisme admet une section (la
diagonale) ; un nerf de Čech augmenté d'un morphisme à section est
contractile. Il en résulte que le complexe augmenté
\[
  0 \to M \to \mathcal{Q}^{0}[M] \to \mathcal{Q}^{1}[M] \to \cdots
\]
est exact, comme complexe de pro-Modules et pour les structures
extrême-droites : \emph{$\mathcal{Q}^{*}[M]$ est une résolution de $M$}.

La page 8 voit aussitôt la difficulté : les pro-Modules ne forment une
catégorie abélienne commode que sous des conditions de type Artin--Rees, et
c'est pour la \emph{limite} $\varprojlim \mathcal{Q}^{*}[M]$ qu'on veut une
résolution. On a la résolution limite sous une hypothèse de
quasi-cohérence sur $M$, et de Mittag-Leffler sur les systèmes
projectifs.\footnote{L'hypothèse est nommée à la page 8, « hypothèse de
quasi-cohérence sur $M$ », au milieu de mots illisibles ; la condition de
Mittag-Leffler, que la page ne nomme pas, est ce qui la rend suffisante :
elle assure que la limite des complexes exacts reste exacte. Nous
l'ajoutons.}

\pagerange{9}{9}
\section*{La cohomologie d'un Module par la résolution canonique (page 9)}

Sous ces conditions, la résolution donne
\[
  H^{*}(X, M) \simeq \mathbb{H}^{*}\bigl(X, \varprojlim
  \mathcal{Q}^{*}[M]\bigr) ,
\]
et la suite spectrale du complexe double,
\[
  E_{2}^{pq} = H^{p}\Bigl(\nu \mapsto H^{q}\bigl(X, \varprojlim
  \mathcal{Q}^{\nu}[M]\bigr)\Bigr) \Longrightarrow H^{p+q}(X, M) .
\]
Si l'on note $\mathcal{Q}^{\nu}_{(n)}$ le quotient d'ordre $n$, de sorte que
$\mathcal{Q}^{\nu}[M] = \varprojlim_n \mathcal{Q}^{\nu}_{(n)} \otimes M$, on
aura, sous des conditions convenables,
\[
  H^{q}\bigl(X, \varprojlim \mathcal{Q}^{\nu}[M]\bigr)
  = \varprojlim_{n} H^{q}\bigl(X, \mathcal{Q}^{\nu}_{(n)} \otimes M\bigr) ,
\]
le terme de droite étant la cohomologie du $n$-ème voisinage de la diagonale
dans $X^{\nu+2}$.\footnote{La page précise : $X^{\nu+2}$, et non $X^{\nu+1}$,
puisque $\mathcal{Q}^{\nu} = \mathcal{P}^{\nu+1}$ ; c'est l'objet de la note
en marge, « Mention $+1$ ! », et d'un « $= H^{q}(X^{[i]}, M)$ » biffé au bout
de la ligne. Les « conditions convenables » ne sont pas précisées ; ce sont
celles qui annulent le $\varprojlim^{1}$.}

Tout cela s'étend aux complexes. Un complexe $M^{\cdot}$ dont les
différentielles sont des \emph{opérateurs différentiels} — un objet de ce que
la page 9 note $D^{+}(X/S)$ — donne, par la fonctorialité de la section
précédente, un complexe double $\mathcal{Q}^{*}[M^{\cdot}]$, et
$\varprojlim \mathcal{Q}^{*}[M^{\cdot}]$ en est une résolution ; d'où
\[
  \mathbb{H}^{*}(X, M^{\cdot}) \simeq
  \mathbb{H}^{*}\bigl(X, \varprojlim \mathcal{Q}^{*}[M^{\cdot}]\bigr) \simeq
  \mathbb{H}^{*}_{\mathrm{strat}}\bigl(X, \mathcal{Q}^{0}[M^{\cdot}]\bigr) .
\]
Le complexe de Modules stratifiés $\mathcal{Q}^{0}[M^{\cdot}]$, dont les
différentielles sont maintenant \emph{linéaires}, s'appelle la
\emph{formalisation} du complexe d'opérateurs différentiels
$M^{\cdot}$.\footnote{« $D^{+}(\mathrm{Mod}(X/S))$ » est à prendre avec
précaution : $\mathrm{Mod}(X/S)$ n'est pas abélienne, et sa catégorie
dérivée n'est pas définie par la recette usuelle. Ce qui a un sens, et ce que
la formule utilise, est la catégorie des complexes d'opérateurs
différentiels, avec pour équivalences celles qui deviennent des
quasi-isomorphismes après formalisation. Le cadre en a été posé plus tard
(Herrera et Lieberman, 1971 ; M.~Saito, 1989, pour les complexes
différentiels et les $\mathcal{D}$-modules induits). Le dernier isomorphisme,
avec la cohomologie stratifiante, anticipe sur la définition de la page 12.}

\pagerange{10}{13}
\section*{Modules stratifiés de cohomologie, et une suite spectrale (pages 10 à 13)}

Soit encore $M^{\cdot}$ un complexe d'opérateurs différentiels. Le complexe
$\mathcal{Q}^{0}[M^{\cdot}]$ est un complexe de pro-Modules stratifiés ; ses
différentielles sont linéaires pour la structure gauche, non pour la
structure extrême-droite. On en prend les faisceaux de cohomologie
\[
  V^{q}(M^{\cdot}) =
  \mathcal{H}^{q}\bigl(\varprojlim \mathcal{Q}^{0}[M^{\cdot}]\bigr) ,
\]
et l'on suppose que leur formation commute aux changements de base
$\mathcal{P}^{0} \to \mathcal{P}^{\nu}$ par les cofaces. C'est une hypothèse
de platitude ; elle est satisfaite, par exemple, quand $X$ est lisse sur
$S$.\footnote{La page 11 est lacunaire ; on y lit une « platitude
infinitésimale de $X/S$ », une platitude de $X^{2}/X \to X$, et « p.\,ex.\
$X/S$ lisse ». Nous ne retenons que l'exemple, qui suffit à la suite.}
Alors :

\begin{enumerate}
\item[1°)] les $V^{q}$ sont, de façon naturelle, des $\mathcal{O}_X$-Modules
  \emph{stratifiés} ;
\item[2°)] dans le complexe double $\varprojlim \mathcal{Q}^{*}[M^{\cdot}]$,
  la première cohomologie, prise suivant les différentielles venues de
  $M^{\cdot}$, est formée des complexes de Čech--Alexander $C^{*}(V^{q})$ des
  stratifications des $V^{q}$ ;
\item[3°)] d'où une suite spectrale.
\end{enumerate}

Définissons, pour un Module stratifié $\underline{V}$ relativement à $S$, sa
\emph{cohomologie stratifiante}
\[
  \mathbb{H}^{p}_{\mathrm{strat}}(X/S, \underline{V}) =
  \mathbb{H}^{p}\bigl(X, C^{*}(\underline{V})\bigr) .
\]
La suite spectrale du 3°) s'écrit alors
\[
  \boxed{\;
  E_{2}^{pq} = \mathbb{H}^{p}_{\mathrm{strat}}\bigl(X/S, V^{q}(M^{\cdot})\bigr)
  \Longrightarrow \mathbb{H}^{p+q}(X, M^{\cdot})
  \;}
\]
\footnote{La page 12 ajoute « vérif.\ à EGA $0_{\mathrm{III}}$ ? » : il
renvoie au formalisme des suites spectrales de complexes doubles d'EGA~0$_{\mathrm{III}}$,
§\,11, pour la vérification. « EGA » est lu sous réserve.}

\emph{Corollaire.} Si $V^{q}(M^{\cdot}) = 0$ pour $q \neq q_0$, on a des
isomorphismes canoniques
\[
  \mathbb{H}^{n}(X, M^{\cdot}) \simeq
  \mathbb{H}^{n-q_0}_{\mathrm{strat}}\bigl(X/S, V^{q_0}(M^{\cdot})\bigr) .
\]
Autrement dit — le cas $q_0 = 0$ est celui qu'il a en vue —
l'hypercohomologie d'un complexe d'opérateurs différentiels qui « résout »
un Module stratifié $\underline{V}$ ne dépend que de $\underline{V}$, et non
du complexe choisi pour le résoudre.

\pagerange{13}{15}
\section*{Le complexe de De Rham et le lemme de Poincaré formel (pages 13 à 15)}

Un Module stratifié $\underline{V}$ sur $X/S$ porte en particulier une
connexion intégrable, et donc un \emph{complexe de De Rham}
\[
  \Omega^{\cdot}(\underline{V}) = \Omega^{\cdot}_{X/S} \otimes_{\mathcal{O}_X}
  \underline{V} ,
\]
complexe d'opérateurs différentiels d'ordre $1$. Lorsque $X$ est lisse sur
$S$ et que $S$ est de caractéristique nulle, ce complexe est une
« résolution différentielle » de $\underline{V}$ au sens de la section
précédente : c'est le \emph{lemme de Poincaré formel}, que la page 13 nomme
en marge.\footnote{La condition sur $S$ est illisible à la page 13 (« et $S$
[\ldots] de [\ldots] ») ; la page 15 la donne : « $X/S$ différentiellement
lisse et $S$ de car.\ nulle ». Elle est nécessaire : en caractéristique $p$,
le complété formel du complexe de De Rham d'une droite n'est pas exact
($t^{p-1}\,dt$ n'a pas de primitive), et c'est précisément ce que
répareront les puissances divisées.} Le corollaire précédent, avec $q_0 = 0$,
donne alors
\[
  H^{*}_{\mathrm{DR}}(X/S, \underline{V}) \simeq
  H^{*}_{\mathrm{strat}}(X/S, \underline{V}) ,
\]
une description de la cohomologie de De Rham où il n'est plus question de
formes différentielles.

Sans hypothèse de lissité ni de caractéristique, il reste un homomorphisme
canonique $\underline{V} \to \mathcal{H}^{0}\bigl(\varprojlim
\mathcal{Q}^{0}[\Omega^{\cdot}(\underline{V})]\bigr)$, c'est-à-dire un
morphisme de complexes de Modules stratifiés
$\underline{V} \to \varprojlim \mathcal{Q}^{0}[\Omega^{\cdot}(\underline{V})]$,
et par la formalisation de la page 9, un homomorphisme canonique
\[
  \boxed{\;
  H^{*}_{\mathrm{strat}}(X/S, \underline{V}) \longrightarrow
  H^{*}_{\mathrm{DR}}(X/S, \underline{V})
  \;}
\]
qui est un isomorphisme lorsque $X$ est lisse sur $S$ et $S$ de
caractéristique nulle.

En caractéristique $p$, il doute que ce soit un isomorphisme, même dans le
cas le plus favorable : $X$ lisse et projectif sur un corps algébriquement
clos de caractéristique $p$, $\underline{V} = \mathcal{O}_X$ avec sa
stratification canonique. Il n'est même pas clair \emph{a priori} que les
$H^{i}_{\mathrm{strat}}$ soient de dimension finie, ni qu'ils s'annulent en
grand degré : « Situation à éclaircir ! »\footnote{La page 15 cite une courbe
elliptique comme exemple de ce doute. Un fait élémentaire l'explique : en
caractéristique $p$, une section horizontale pour une stratification est
annulée par tous les opérateurs différentiels, donc est une puissance
$p^{n}$-ième pour tout $n$ ; la cohomologie stratifiante de $\mathcal{O}_X$
ne voit donc presque rien de la cohomologie de De Rham. Pour ce qu'on sait
aujourd'hui de la cohomologie du site infinitésimal en caractéristique $p$,
voir l'article d'A.~Ogus, « The cohomology of the infinitesimal site ».}

\emph{Remarque.} L'intérêt de ce lien est que la cohomologie stratifiante
s'interprète, à son tour, comme la cohomologie d'un « site stratifiant », et
qu'une variante convenable de ce site pourrait conduire à la cohomologie
$p$-adique cherchée en caractéristique $p$. Une note en marge, dans un cadre,
y répond : « C'est sûrement faux ! », avec un renvoi à des calculs sur les
variétés abéliennes.\footnote{La première ligne de la note est au crayon, les
suivantes à l'encre bleue : elle a été ajoutée après coup, sans doute après
le calcul de la page 16. « $p$-adique » et « abélien » sont lus sous
réserve.}

\pagerange{16}{18}
\section*{Commentaires sur une cohomologie $p$-adique en caractéristique $p$ (pages 16 à 18)}

Soit $X_0$ un schéma sur $k$, et $S = \operatorname{Spec} W$. Le premier espoir
était que le $W$-module gradué
\[
  H^{*}\bigl((X_0/S)_{\mathrm{inf}}, \mathcal{O}\bigr)
\]
soit « le » bon candidat. Si $X_0$ est lisse et se relève en un $X$ lisse sur
$W$, de réductions $X_n$ sur $W_n$, on a, « d'après l'exposé précédent » — sans doute l'invariance des pages 64
et 65 —
\[
  H^{*}\bigl((X_0/S)_{\mathrm{inf}}, \mathcal{O}\bigr) = \varprojlim_{n}
  H^{*}\bigl((X_n/S_n)_{\mathrm{inf}}, \mathcal{O}\bigr) ,
\]
et l'on espérait trouver à la limite la même chose que $\varprojlim_n
H^{*}_{\mathrm{DR}}(X_n/W_n) \simeq H^{*}_{\mathrm{DR}}(X/W)$.

« C'est faux malheureusement. » Si $X_0$ est une variété abélienne
\emph{ordinaire} — formellement, un tore — on trouve
\[
  H^{i}\bigl((X_0/S)_{\mathrm{inf}}, \mathcal{O}\bigr) = 0
  \quad \text{pour } i > \dim X_0 ,
\]
alors que la limite des cohomologies de De Rham est excellente. Dans le cas
affine, on voit d'où vient l'ennui : le passage à la limite projective ne
commute pas à la localisation $W \to W[1/p] = K$.\footnote{Le défaut est
celui du site sans puissances divisées ; c'est en adjoignant des puissances
divisées aux idéaux des épaississements que la théorie cristalline de
Berthelot (1974) obtient, pour $X_0$ propre et lisse relevable,
$H^{*}_{\mathrm{cris}}(X_0/W) \simeq H^{*}_{\mathrm{DR}}(X/W)$. Le dossier ne
contient pas cette idée.}

\emph{Une proposition.} Remplacer le site par un site plus gros, formé
d'épaississements \emph{pas forcément infinitésimaux} : les objets sont les
triples $(U_0, \mathcal{W}, i)$ où $U_0$ est un ouvert affine de $X_0$,
$\mathcal{W} = \operatorname{Spf} B$ un schéma formel affine complet sur $W$,
et $i : U_0 \to \mathcal{W}$ un morphisme qui identifie $U_0$ à un
sous-schéma fermé de la fibre spéciale $\mathcal{W} \times_S S_0$, défini
par un idéal nilpotent — en termes d'anneaux, un $W$-homomorphisme surjectif
$B \to A_0$ dont le noyau, réduit modulo $p$, est nilpotent. Topologie : une
variante convenable de la topologie de Zariski, les recouvrements se
lisant sur les traces dans $X_0$.\footnote{La page 17 écrit « isomorphisme
nilpotent » ; nous l'entendons comme immersion fermée définie par un idéal
nilpotent, ce que confirme la parenthèse sur $B/pB \to A_0$. Le site proposé
est à rapprocher, sans s'y identifier, du site convergent d'Ogus (« $F$-isocrystals
and de Rham cohomology II — Convergent isocrystals », 1984), dont les objets,
les « enlargements », sont aussi des schémas formels $p$-adiques munis d'une
application de leur fibre spéciale réduite vers $X_0$, et de la cohomologie
rigide de Berthelot. Rien n'indique que ces travaux s'appuient sur cette
page.}

\emph{Tests décisifs.}
a) Si $X_0$ est lisse et affine sur $k$, on doit retrouver la cohomologie de
Washnitzer--Monsky.
b) Si $X_0$ est lisse et propre et se relève en $X$ lisse et propre, on doit
retrouver, au moins modulo torsion, $H^{*}_{\mathrm{DR}}(X/W)$.

« Il s'agit d'une première [tentative]. » Même si elle réussit, il faudra
beaucoup pour en faire la théorie cohomologique désirée.\footnote{La fin de
la page 18 est presque entièrement illisible ; « [tentative] » est notre
restitution d'un mot biffé, « W.M. » (au test a) et « accomplir » sont lus
sous réserve. Le dossier ne revient pas sur ces tests.}

\pagerange{20}{21}
\section*{Le « yoga » de la cohomologie de De Rham (pages 20 à 26)}

L'exposé de novembre 1966 s'annonce, dans un cadre ajouté en tête : « faire
sentir un complexe de problèmes gravitant autour de la cohomologie de De Rham
des variétés algébriques, et proposer un point de vue qui me semble permettre
de les grouper et de les étudier. D'où un (gros) programme de travail
proposé. »

\emph{Variétés différentiables.} Le théorème de De Rham s'explique par les
faisceaux : $\Omega^{\cdot}_X$ est une résolution de $\mathbb{C}_X$ (lemme de
Poincaré), et les $\Omega^{i}_X$ sont fins, donc acycliques ($X$ séparé,
paracompact).

\emph{Variétés analytiques complexes.} Le lemme de Poincaré holomorphe vaut
encore, mais les $\Omega^{i}_X$ ne sont plus acycliques, sauf sur une variété
de Stein. En général il faut l'hypercohomologie :
\[
  H^{*}(X, \mathbb{C}) \simeq \mathbb{H}^{*}(X, \Omega^{\cdot}_X) ,
  \qquad
  (*)\quad E_{1}^{pq} = H^{q}(X, \Omega^{p}_X) \Longrightarrow
  \mathbb{H}^{p+q}(X, \Omega^{\cdot}_X) ,
\]
la suite spectrale de Hodge vers De Rham.\footnote{Un « 12 » de lecture
douteuse surmonte la première égalité ; nous ne l'interprétons pas.}

\emph{Variétés algébriques non singulières} sur $\mathbb{C}$, avec la
topologie de Zariski. On a toujours un homomorphisme
\[
  \mathbb{H}^{*}(X, \Omega^{\cdot}_X) \longrightarrow
  \mathbb{H}^{*}(X^{\mathrm{an}}, \Omega^{\cdot}_{X^{\mathrm{an}}})
  \simeq H^{*}(X^{\mathrm{an}}, \mathbb{C}) .
\]

\emph{Théorème de comparaison.} Pour $X$ lisse et localement de type fini sur
$\mathbb{C}$, c'est un isomorphisme. Si $X$ est affine, cela revient à dire
que $H^{*}(X^{\mathrm{an}}, \mathbb{C})$ se calcule par le complexe des formes
différentielles \emph{algébriques} globales.\footnote{La parenthèse de la
page 21 dit seulement « $X$ schéma loc.\ de type fini sur $\mathbb{C}$ » ;
la lissité, que le contexte (« Variétés algébriques (non singulières) »,
page 20) suppose, est nécessaire : la page 23 donne elle-même le
contre-exemple d'une courbe nodale. La démonstration renvoie à un « papier
bleu », sans doute la lettre à Atiyah publiée sous le titre « On the de Rham
cohomology of algebraic varieties », \emph{Publ.\ Math.\ IHÉS} 29 (1966) —
identification de notre part.} Le cas général se ramène aux cas
particuliers par la suite spectrale d'un recouvrement. Si $X$ est propre,
cela résulte des deux suites spectrales $(*)$, algébrique et analytique, et
de GAGA ; le cas général demande un moyen plus puissant que GAGA, la
résolution des singularités.

\pagerange{21}{22}
\subsection*{Conséquences « métaphysiques », et critique (pages 21 et 22)}

La cohomologie complexe, définie par voie transcendante, admet une
description purement algébrique, qui a un sens pour des schémas lisses sur
un corps quelconque ; elle est de plus \emph{filtrée}, le gradué associé
étant le terme $E_\infty$ de la suite spectrale $(*)$ — la filtration de
Hodge.

\emph{Critique.} En caractéristique nulle, pour les variétés complètes et
lisses, elle donne des invariants excellents, et les bons nombres de Betti.
En caractéristique $p$, elle n'est pas raisonnable :
\begin{enumerate}
\item[1°)] si $X$ n'est pas propre, on trouve des espaces de dimension
  infinie : pour $X$ lisse sur un corps parfait, $H^{0}_{\mathrm{DR}}(X) =
  \Gamma(X, \mathcal{O}_X)^{p}$, l'anneau des puissances
  $p$-ièmes ;\footnote{Le noyau de $d$ sur $\mathcal{O}_X$ est
  $\mathcal{O}_X^{p}$ pour $X$ lisse sur un corps parfait (Cartier) ; la page
  22 n'énonce pas l'hypothèse, que nous ajoutons.}
\item[2°)] même si $X$ est propre, on ne trouve pas en général les bons
  nombres de Betti : on a $\dim H^{1}_{\mathrm{DR}}(X) \geq 2 \dim
  \mathrm{Pic}_X$, avec inégalité stricte dans certains exemples, dus
  « moralement » à de la $p$-torsion ;\footnote{$\dim \mathrm{Pic}_X$ est la
  dimension de la variété de Picard, et $2\dim\mathrm{Pic}_X$ le premier
  nombre de Betti $\ell$-adique. Les exemples où l'inégalité est stricte
  sont liés à la non-réduction du schéma de Picard, dont la surface d'Igusa
  (1955) est un exemple.}
\item[3°)] même là où les dimensions sont les bonnes, les coefficients sont
  de caractéristique $p$ : la formule de Lefschetz n'y donne les nombres de
  points fixes que modulo $p$, et ne se prête pas aux conjectures de Weil —
  défaut déjà signalé par Serre.
\end{enumerate}

\pagerange{23}{26}
\subsection*{Singularités, suites de Leray, dualité (pages 23 à 26)}

Deux autres critiques valent même sur $\mathbb{C}$. D'abord, le point de vue
est trop exclusivement lié aux variétés \emph{non singulières} : pour la
droite affine ou projective avec un point double ordinaire, $\dim
H^{1}(X, \Omega^{\cdot}_X)$ n'est pas le premier nombre de Betti. Or on ne
peut éviter les variétés singulières, ne serait-ce que comme fibres d'un
morphisme $f : X \to Y$ ; on voudrait une construction purement algébrique
de la suite spectrale de Leray de $f^{\mathrm{an}}$,
\[
  E_{2}^{pq} = H^{p}\bigl(Y^{\mathrm{an}}, R^{q}f^{\mathrm{an}}_{*}
  \mathbb{C}\bigr) \Longrightarrow H^{p+q}(X^{\mathrm{an}}, \mathbb{C}) ,
\]
ce qui exige des coefficients plus généraux que les seuls complexes de De
Rham, et sans doute des complexes d'opérateurs différentiels plus généraux.
Ensuite, la dualité : pour $X$ lisse et propre on a la dualité habituelle,
Gysin, la classe d'un cycle, comme dans son exposé Bourbaki sur la
« cohomologie de Hodge » ; mais on voudrait un formalisme $f_{!}$, $f^{!}$
pour des morphismes séparés de type fini quelconques, du type du Séminaire
Hartshorne (\emph{Residues and Duality}, 1966), qui donnerait aussi
l'homologie singulière, et qui s'étende des complexes à différentielles
linéaires aux complexes d'opérateurs différentiels.

La question posée est donc celle de la bonne \emph{catégorie de
coefficients}, qui englobe à la fois les complexes cohérents
« linéaires » et les complexes de type De Rham.\footnote{C'est la question à
laquelle répondront, en caractéristique nulle, la théorie des
$\mathcal{D}$-modules (Kashiwara, Bernstein, Mebkhout, dans les années 1970)
et son formalisme des six opérations ; ces noms sont postérieurs au dossier.}

\pagerange{26}{32}
\section*{Tour d'horizon $\ell$-adique (pages 26 à 32)}

Pourquoi tant d'efforts pour algébriser De Rham, quand on dispose des
invariants $\ell$-adiques ? Pour $X$ algébrique sur un corps algébriquement
clos $k$ et $\ell \neq \mathrm{car}\,k$,
\[
  H^{i}(X, \mathbb{Z}_\ell) = \varprojlim_{\nu}
  H^{i}(X_{\mathrm{\acute{e}t}}, \mathbb{Z}/\ell^{\nu}\mathbb{Z}) ,
\]
et sur $\mathbb{C}$ le théorème de comparaison donne des isomorphismes
canoniques, fonctoriels en $X$,
\[
  \boxed{\;H^{i}(X, \mathbb{Z}_\ell) \simeq
  H^{i}(X^{\mathrm{an}}, \mathbb{Z}) \otimes_{\mathbb{Z}} \mathbb{Z}_\ell\;}
\]
La connaissance de \emph{tous} les $H^{i}(X, \mathbb{Z}_\ell)$ équivaut donc,
du point de vue additif, à celle de la cohomologie entière transcendante :
pour un $\ell$ donné, au rang $b_i(X)$ et à la composante $\ell$-primaire de
la torsion. On en tire, en caractéristique nulle et par voie transcendante,
que les $H^{i}(X, \mathbb{Z}_\ell)$ sont de type fini, de rang indépendant de
$\ell$, sans torsion pour presque tout $\ell$, et que le polynôme
caractéristique d'un endomorphisme $f$ de $X$ sur $H^{i}(X, \mathbb{Z}_\ell)$
est à coefficients entiers indépendants de $\ell$.

En caractéristique $p$, rien de tel. La finitude n'est connue que pour $X$
propre ou $\dim X \leq 2$ ; pour $X$ projectif et lisse, on ignore si le
rang de $H^{2}$ est indépendant de $\ell$, et à plus forte raison si les
polynômes caractéristiques le sont — ce serait, dit-il, le pas le plus
important vers les conjectures de Weil.\footnote{L'état décrit est celui de
1966. La finitude en général a été établie par Deligne
(« Théorèmes de finitude en cohomologie $\ell$-adique », SGA~4½, 1977), sans
résolution des singularités ; l'indépendance de $\ell$ pour le Frobenius
d'une variété projective et lisse sur un corps fini résulte de la
démonstration des conjectures de Weil (Deligne, 1974). Pour un endomorphisme
quelconque, la question reste liée aux conjectures standard.} Et surtout —
« c'est là le point que je désire faire » — même si l'on montre que ces
coefficients sont rationnels (on sait ramener la question, pour $X$ projectif
lisse, à un problème de cycles algébriques du type de Lefschetz), on ne pourra
conclure qu'ils sont dans $\mathbb{Z}[1/p]$, à cause de la restriction
$\ell \neq p$. \emph{Il manque une théorie cohomologique $p$-adique pour les
variétés en caractéristique $p$} ; sans elle on ne comprendra pas non plus
les phénomènes de $p$-torsion, dont l'exposé de Tate sur le groupe de Brauer
donne des exemples.\footnote{La note oblique de la marge de la page 31, qui
porte sur le groupe de Brauer sur un corps fini, est presque entièrement
illisible.}

Que serait une telle théorie ? Un foncteur $X \mapsto H^{\cdot}(X)$ des
variétés projectives lisses sur $k$ vers les modules de type fini sur
$\mathbb{Z}_p$ ou sur $\mathbb{Q}_p$, avec Künneth, dualité, et les bons
nombres de Betti. Un argument de Serre montre qu'un tel foncteur n'existe
pas, déjà pour le $H^{1}$ et en se bornant à des variétés de dimension
$\leq 2$.\footnote{La page 31 écrit « VA de $\dim \leq 2$ », que l'on peut
lire variétés algébriques ou variétés abéliennes. L'argument de Serre, sous
sa forme classique, porte sur une courbe elliptique supersingulière $E$ :
$\mathrm{End}(E) \otimes \mathbb{Q}$ est une algèbre de quaternions ramifiée
en $p$, qui n'a pas de représentation de dimension $2$ sur $\mathbb{Q}_p$,
alors qu'elle devrait agir sur $H^{1}(E)$.} L'anneau de coefficients naturel
est l'anneau $W(k)$ des vecteurs de Witt, de valuation discrète pour $k$
parfait, et dont le corps des fractions $K \supset \mathbb{Q}_p$ est de
caractéristique \emph{nulle}.

\pagerange{32}{34}
\section*{Essais de cohomologie $p$-adique (pages 32 à 34)}

\emph{Cohomologie plate à coefficients radiciels.} Prendre pour coefficients
des schémas en groupes finis radiciels ($\mu_p$, $\alpha_p$, \ldots), avec la
topologie plate. C'est magnifique en dimension $1$ — dualité en cohomologie
plate pour les courbes sur un corps éventuellement fini, que Raynaud devrait
tirer au clair. Mais sans espoir en dimension $n \geq 2$ : Artin a remarqué
que ces groupes de cohomologie s'annulent en degré $> n + 1$, au lieu de
$2n$. Il faut donc une \emph{définition de la cohomologie $p$-adique
indépendante de la cohomologie plate}, quitte à étudier ensuite ses
relations, sans doute fort intéressantes, avec celle-ci.

\pagerange{34}{37}
\section*{Washnitzer--Monsky (pages 34 à 46)}

\emph{L'idée, simplifiée à l'extrême.} Soit $X_0$ lisse sur $k$, et
supposons qu'il se relève en un $X$ propre et lisse sur $S =
\operatorname{Spec} W$ :

\begin{tikzcd}
X_0 \arrow[r, hook] \arrow[d] & X \arrow[d] \\
\operatorname{Spec} k \arrow[r, hook] & \operatorname{Spec} W
\end{tikzcd}

On pose
\[
  \boxed{\;H^{\cdot}_{\mathrm{DR}}(X/S) = \mathbb{H}^{*}(X, \Omega^{\cdot}_{X/S})\;}
\]
module de type fini sur $W$.\footnote{La page 34 dessine ce carré en marge,
avec des traits verticaux sans flèche ; nous mettons les flèches structurales,
qui sont ce que les traits signifient.} L'idée de Washnitzer et Monsky
semble pouvoir s'exprimer en disant que ce module, à isomorphisme canonique
près, \emph{ne dépend pas du relèvement} $X$, et serait donc un invariant de
$X_0$ — par définition la cohomologie $p$-adique cherchée.

La reprise du 15 novembre (page 35) rappelle le cadre et ajoute que
$H_{\mathrm{DR}}(X/S)$ et les $Rf_{*}\Omega^{\cdot}_{X/S}$ ont un sens pour
tout préschéma relatif, et commutent au changement de base si $f$ est
lisse.\footnote{La note encadrée en bas de la page 35 est en partie
illisible ; « comp.\ avec chgt de base si $f$ lisse » est la partie lue, à
entendre pour $f$ propre et lisse.} Puis trois remarques :
\begin{enumerate}
\item[a)] la construction marcherait encore si l'on ne peut relever $X_0$
  qu'en un schéma \emph{formel} lisse sur $S$ ;
\item[b)] le texte publié de Washnitzer--Monsky ne traite que
  $H_{\mathrm{DR}} \otimes_W K$, et seulement pour $X_0$ affine, relevé non en
  un schéma mais en quelque chose qui ressemble à une variété
  rigide-analytique à la Tate ; ce sont des invariants \emph{locaux}.
  D'après Washnitzer, il serait aisé d'obtenir des invariants globaux, sur
  $W$, en travaillant sur le \emph{site} des relèvements des ouverts
  affines ;
\item[c)] quand $X_0$ admet un relèvement propre et lisse, la construction
  donne les bons nombres de Betti : ceux de la fibre générique, par le
  théorème de comparaison.\footnote{Ce que la page 36 appelle « l'invariance
  de W.M. » pour un relèvement propre et lisse sur $W$ est aujourd'hui vrai
  et plus fort qu'il ne le pensait : $H^{*}_{\mathrm{DR}}(X/W) \simeq
  H^{*}_{\mathrm{cris}}(X_0/W)$ (Berthelot, 1974), sans tensoriser par $K$.
  Mais c'est la cohomologie cristalline à puissances divisées qui le
  démontre, non la construction de Washnitzer--Monsky. La note en marge de la
  page 36, « Some properties of formal schemes 63/64 », est une référence que
  nous n'identifions pas.}
\end{enumerate}

\pagerange{38}{40}
\subsection*{Pourquoi le complété formel ne suffit pas (pages 38 à 40)}

Pour $X_0$ affine et lisse, on sait relever $X_0$ en un schéma formel lisse
$\mathfrak{X}$ sur $W$, unique à isomorphisme non unique près (SGA~1,
exposé~III). La première idée serait de prendre
$\mathbb{H}^{*}(\mathfrak{X}, \Omega^{\cdot}_{\mathfrak{X}/S}) =
H^{*}\bigl(\Gamma(\mathfrak{X}, \Omega^{\cdot}_{\mathfrak{X}/S})\bigr)$. C'est
faux, déjà pour la droite affine : $X_0 = \operatorname{Spec} k[t]$,
$\mathfrak{X} = \operatorname{Spf} W\{t\}$, où $W\{t\} = \varprojlim_n
W_n[t]$ est l'anneau des séries restreintes, à coefficients tendant vers $0$.
On voudrait $\mathbb{H}^{1} = 0$, ou au moins de torsion. Or
\[
  \mathbb{H}^{1}(\mathfrak{X}, \Omega^{\cdot}_{\mathfrak{X}/S})
  \simeq W\{t\} \big/ \{ f' : f \in W\{t\} \} ,
\]
et $\sum a_n t^{n}$ est une dérivée dans $W\{t\}$ si et seulement si $\sum
a_n t^{n+1}/(n+1)$ est encore restreinte, c'est-à-dire
\[
  v_p(a_n) - v_p(n+1) \longrightarrow +\infty .
\]
Comme $v_p(n+1)$ prend des valeurs arbitrairement grandes, la seule condition
$v_p(a_n) \to +\infty$ ne suffit pas, et $\mathbb{H}^{1}$ n'est pas de
torsion : il est de rang infini sur $W$.

La condition $|a_n|_p / |n+1|_p \to 0$ est une condition arithmétique
« assez biscornue ». Washnitzer et Monsky la remplacent par une condition plus
forte et plus maniable :
\[
  |a_n|_p \leq C \rho^{-n} \quad \text{pour un } \rho > 1
  \text{ et une constante } C
\]
(dépendant de la série), c'est-à-dire la convergence sur un disque de rayon
$\rho > 1$, un peu plus grand que le disque unité. C'est l'anneau
$W^{\dagger}\{t\}$ des séries \emph{surconvergentes} ; on définit de même
$W^{\dagger}\{t_1, \ldots, t_n\}$, et les algèbres de Washnitzer--Monsky,
quotients de celles-ci par des idéaux quelconques.\footnote{La page 40 écrit
$|a_n| = O(1)\,\rho^{n}$, $\rho > 1$ : c'est une majoration de croissance, qui
ne donne aucune convergence ; la condition voulue, celle de son propre
commentaire (« disque de rayon $\rho > 1$ »), est $|a_n| = O(\rho^{-n})$. Elle
donne bien $|a_n/(n+1)| \leq C (n+1) \rho^{-n} = O(\rho'^{-n})$ pour tout
$1 < \rho' < \rho$, de sorte que toute série surconvergente a une primitive
surconvergente \emph{à coefficients dans $K$} : sur $W$, $H^{1}$ de la droite
surconvergente est de torsion (non bornée), et nul après $\otimes K$. La même
page dit « $K$ corps des fractions de $K$ », lapsus pour « de $W$ ».} Les
idéaux sont-ils automatiquement fermés ? La question reste posée en tête de
la page 41, et le dossier ne la tranche pas.\footnote{On sait depuis que les
algèbres de Washnitzer--Monsky sont noethériennes (Fulton, « A note on weakly
complete algebras », 1969) ; pour les fondements, voir M.~van der Put,
« The cohomology of Monsky and Washnitzer », \emph{Mémoires SMF} 23 (1986).}

\pagerange{41}{43}
\subsection*{Le théorème de relèvement et l'invariance (pages 41 à 43)}

Il semble vrai, écrit-il, que :
\begin{enumerate}
\item[a)] si $A_0$ est une $k$-algèbre de type fini lisse, il existe une
  algèbre de Washnitzer--Monsky $A$ sur $W$, plate (et formellement lisse)
  sur $W$, et un isomorphisme $\varphi : A \otimes_W k \simeq A_0$ ;
\item[b)] un tel couple $(A, \varphi)$ est unique à isomorphisme non unique
  près ;
\item[c)] on définit $\Omega^{1}(A/W)$ par la propriété universelle pour les
  dérivations de $A$ dans les $A$-modules de type fini, et
  $\Omega^{\cdot}(A/W) = \bigwedge^{\cdot}_{A} \Omega^{1}(A/W)$, complexe de
  $W$-modules ;\footnote{La page 41 écrit $\bigwedge^{*}(A/W)$, en omettant
  $\Omega^{1}$, et dit « complexe de $W$-algèbres » ; c'est une algèbre
  différentielle graduée sur $W$.}
\item[d)] un homomorphisme $u : A \to A'$ d'algèbres de
  Washnitzer--Monsky lisses induit
  \[
    u_{*} : H^{\cdot}\bigl(\Omega^{\cdot}(A/W)\bigr) \otimes_W K
    \longrightarrow H^{\cdot}\bigl(\Omega^{\cdot}(A'/W)\bigr) \otimes_W K ,
  \]
  qui ne dépend que de l'homomorphisme réduit $u_0 : A_0 \to A'_0$.
\end{enumerate}
\emph{Corollaire.} Pour $A$ un relèvement de Washnitzer--Monsky de $A_0$
lisse, $H^{\cdot}\bigl(\Omega^{\cdot}(A/W)\bigr) \otimes_W K$ ne dépend, à
isomorphisme canonique près, que de $A_0$.\footnote{Deux écarts avec la
page. Le sens de la flèche : la page 41 écrit $u^{*}$ de $A'$ vers $A$, d'un
simple trait sans pointe ; un homomorphisme d'anneaux $A \to A'$ induit une
application de $\Omega^{\cdot}(A)$ vers $\Omega^{\cdot}(A')$. Et le
$\otimes K$ : la page énonce d) et le corollaire sur $W$, comme un théorème
qu'« il semble » vrai, mais la page 42 précise que Washnitzer--Monsky ne
prouvent l'invariance que modulo torsion. Elle est fausse sur $W$ en
général : l'homotopie entre deux relèvements de $u_0$ fait intervenir des
divisions par des factorielles, et la page 40 vient de montrer que
$H^{1}$ de la droite surconvergente a de la torsion non bornée. L'existence
en a) est due à Elkik (1973) pour les relèvements lisses, l'unicité en b) à
Washnitzer et Monsky.}

Washnitzer et Monsky ne prouvent ces résultats que pour des $A_0$
« spéciales », dont il y a assez. Pour $X_0$ lisse quelconque, ils en
tirent des faisceaux de cohomologie $\mathcal{H}^{*}_{X_0}$ sur $X_0$, de
$W$-modules, mais pas encore d'invariants globaux. Récemment, ils ont
utilisé un \emph{site} — celui des relèvements de Washnitzer--Monsky des
ouverts affines spéciaux, pour la topologie de Zariski — sur lequel les
$\Omega^{\cdot}_{A/W}$ forment un complexe de faisceaux dont on prend
l'hypercohomologie ; ils trouvent ainsi des invariants globaux, qui redonnent
dans le cas relevable le $\mathbb{H}^{*}(\mathfrak{X},
\Omega^{\cdot}_{\mathfrak{X}/S})$ de la page 38.\footnote{« vérifient (?) » et
« ils retrouvent bien » sont des ajouts serrés en interligne, dont l'ordre
est reconstitué par la transcription.}

\pagerange{43}{46}
\subsection*{Critique de Washnitzer--Monsky (pages 43 à 46)}

C'est une idée nouvelle importante : elle relie la cohomologie de De Rham en
caractéristique nulle à la cohomologie $p$-adique hypothétique, et évite une
des critiques principales faites à De Rham naïf, puisque ses coefficients
sont dans $W(k)$, de corps des fractions de caractéristique nulle. Si elle
vit vraiment sur $W$, et pas seulement sur $K$, elle fournit en outre une
notion de \emph{torsion $p$-adique}, peut-être la bonne — est-ce celle de la
fibre générique ? géométrique ? C'est la question, « extrêmement importante
et profonde », des relations entre cohomologie de De Rham et cohomologie
$p$-adique dans la spécialisation de la caractéristique $0$ à la
caractéristique $p$, que Tate a envisagée l'année précédente dans un autre
langage.\footnote{Le lieu où Tate l'a envisagée est illisible : on lit deux
capitales entourées, « HH » ou « IH ». C'est aujourd'hui le domaine de la
théorie de Hodge $p$-adique ; la page renvoie aussi, pour son importance dans
l'étude des familles de cycles algébriques, à un texte à paraître dont le
titre est illisible.}

Restent des critiques graves. La théorie est fournie par les formes
différentielles, donc pratiquement limitée aux $X_0$ lisses ; et en ne
travaillant qu'avec les complexes de De Rham, il est sans espoir qu'elle
fournisse des suites spectrales de Leray pour un morphisme $f_0 : X_0 \to
Y_0$, si précieuses pour les dévissages. Sans doute en conséquence,
Washnitzer et Monsky n'ont pas pu prouver les propriétés formelles les plus
élémentaires : que leurs $H^{i}(X_0)$ soient de type fini, donc ni la
dualité ni Künneth.\footnote{La finitude de la cohomologie de
Washnitzer--Monsky, tensorisée par $K$, n'a été démontrée qu'en 1997, dans le
cadre de la cohomologie rigide (Berthelot, « Finitude et pureté
cohomologique en cohomologie rigide », \emph{Invent.\ Math.} 128 ; Mebkhout,
la même année, par une autre voie).} Ils envisagent cependant d'exprimer la
fonction $\zeta$ d'un schéma projectif et lisse sur un corps fini par une
formule de Lefschetz--Weil, Frobenius opérant sur leurs $H^{i}$ — un
formalisme proche de celui de Dwork, « que je ne connais pas ».\footnote{C'est
le contenu de Monsky, « Formal cohomology III : Fixed point theorems »,
\emph{Annals of Math.} 1971.}

\pagerange{46}{49}
\section*{Washnitzer--Monsky et Gauss--Manin (pages 46 à 52)}

L'invariance affirmée par Washnitzer--Monsky surprend d'abord. Supposons que
$X_0$ varie beaucoup : qu'il soit la fibre, en un point $t_0$ au-dessus du
point fermé $s_0$ de $S$, d'un schéma $Y$ propre et lisse sur un $M$ de type
fini sur $S$,

\begin{tikzcd}[column sep=large]
Y \arrow[r, "f"] & M \arrow[r, "h"] & S = \operatorname{Spec} W
\end{tikzcd}

Chaque section $g$ de $M$ sur $S$ passant par $t_0$ donne un relèvement
$X_g = g^{*}Y$ de $X_0$, propre et lisse sur $S$, de morphisme structural
$f_g$. Posons $\mathcal{K}^{\cdot} = Rf_{*}(\Omega^{\cdot}_{Y/M})$ et
$\mathcal{H}^{i} = \mathcal{H}^{i}(\mathcal{K}^{\cdot}) =
R^{i}f_{*}(\Omega^{\cdot}_{Y/M})$. Le changement de base, au sens des
catégories dérivées, donne
\[
  Rf_{g*}(\Omega^{\cdot}_{X_g/S}) \simeq g^{*}\mathcal{K}^{\cdot} ,
  \qquad
  R^{i}f_{g*}(\Omega^{\cdot}_{X_g/S}) \simeq g^{*}\mathcal{H}^{i}
  \ \text{(modulo torsion)} .
\]
Si l'invariance est vraie, les $g^{*}\mathcal{K}^{\cdot}$, pour les
diverses sections $g$ passant par $t_0$, sont canoniquement isomorphes entre
eux, et les $g^{*}\mathcal{H}^{i}$ aussi, modulo torsion.\footnote{Les
exposants des pages 47 et 48, un petit signe entre parenthèses lu « $(\Sigma)$ »
sous réserve, sont omis ; ils n'interviennent pas dans l'argument.} Un tel
système transitif d'isomorphismes est une \emph{structure} sur
$\mathcal{K}^{\cdot}$, ou sur les $\mathcal{H}^{i}$, relativement à $S$ —
apparentée à une donnée de descente : pour un $M$-module $\underline{E}$, un
isomorphisme $\underline{E} \simeq h^{*}\underline{E}_0$ en fournirait un.
Les structures plus fortes qui en fournissent encore un sont les
\emph{stratifications relativement à $S$}. Suivent, page 49, quatre titres
sans texte : définition générale d'une stratification et relation avec les
connexions ; système transitif d'isomorphismes qu'elle fournit sur une base
locale complète ; la connexion de Gauss--Manin.

\pagerange{50}{52}
\subsection*{La connexion de Gauss--Manin (pages 50 à 52)}

En caractéristique nulle, la construction de Manin donne une connexion
intégrable absolue sur la cohomologie de De Rham d'une famille, sans le
langage géométrique ; pour le $H^{1}$ d'une famille de variétés propres et
lisses, Manin montre de façon compliquée que la partie
$\mathbb{H}(X, \Omega^{\cdot}_{X/k})$ est stable par les dérivations
$\theta_X$.\footnote{La page 50 commence au milieu d'un développement, sans
titre : l'exposé de Manin annoncé au point 4 de la page 49 ne figure pas
sur les feuillets. Deux lettres de la page sont incertaines ($k$ ou $h$ sous
$\Omega$, $K$ dans $H_{\mathrm{DR}}(K/k)$). La référence est Manin,
« Algebraic curves over fields with differentiation », 1958.}

En fait, des arguments directs, reposant sur la formule d'homotopie de
Cartan
\[
  d \circ i_\xi + i_\xi \circ d = \theta_\xi ,
\]
($i_\xi$ le produit intérieur, $\theta_\xi$ la dérivée de Lie) donnent, pour
un morphisme $f : X \to S$ :
\begin{enumerate}
\item[(i)] Si une dérivation $\xi$ de $\mathcal{O}_S$ se relève localement en
  une dérivation $\overline{\xi}$ de $\mathcal{O}_X$ — c'est toujours le cas si
  $X$ est formellement lisse et affine sur $S$ —, l'action de
  $(\overline{\xi}, \xi)$ sur $f_{*}(\Omega^{\cdot}_{X/S})$ ne dépend du choix
  de $\overline{\xi}$ qu'à homotopie près. Les dérivations de $S$ opèrent donc
  sur les $R^{i}f_{*}(\Omega^{\cdot}_{X/S})$, avec les lois de linéarité et
  de crochet ; si $S$ est lisse sur un $T$, c'est une connexion intégrable
  relative à $T$.
\item[(ii)] Par le même principe, un peu plus sophistiqué (renvoyé à un
  « appendice » qui n'est pas dans le dossier), pour $f$ lisse et propre,
  $Rf_{*}(\Omega^{\cdot}_{X/S})$ est muni d'une connexion absolue canonique au
  sens des catégories dérivées. Cela ne munit pas individuellement les
  $R^{i}f_{*}$ de connexions ; c'est pourtant le cas si $S$ est de
  caractéristique nulle, car alors la formation des $R^{i}f_{*}$ commute au
  changement de base. La compatibilité avec la construction transcendante est
  triviale.\footnote{La page 51 dit « lisse propre et qu.-sép. ». Une
  construction de ce type, par filtration du complexe de De Rham absolu, est
  publiée par Katz et Oda, « On the differentiation of De Rham cohomology
  classes with respect to parameters », 1968.}
\end{enumerate}

\emph{Remarque.} On pourrait espérer que ces connexions proviennent de
\emph{stratifications} canoniques, ce qui expliquerait \emph{a priori}
l'invariance de Washnitzer--Monsky. C'est vrai en caractéristique nulle ;
c'est faux en caractéristique $p$, par exemple pour les familles de courbes
elliptiques. Il s'agit donc d'une relation « heuristique et sentimentale »
entre deux idées, dont le lien exact reste à élucider.\footnote{En
caractéristique $p$, une stratification force la $p$-courbure de la connexion
associée à s'annuler, et celle de Gauss--Manin pour une famille de courbes
elliptiques non isotriviale ne s'annule pas (Katz, « Nilpotent connections and
the monodromy theorem », 1970). Le lien exact viendra des puissances
divisées : la cohomologie cristalline d'une famille propre et lisse est un
\emph{cristal}, c'est-à-dire porte une stratification à puissances divisées,
dont la connexion sous-jacente est Gauss--Manin.}

Pour résumer, on est devant un puzzle : reconstituer une théorie d'ensemble
qui combine (a) une théorie cohomologique des complexes d'opérateurs
différentiels, en vue notamment du complexe de De Rham ; (b) une étude
systématique des structures stratifiées ; (c) la théorie de
Washnitzer--Monsky.

\pagerange{53}{55}
\section*{Topos stratifiant et topos infinitésimal (pages 53 à 55)}

\emph{4.1. Ce qu'est une stratification.} Soit $\underline{F}$ un Module sur
$X$ muni d'une stratification relative à $S$. Pour tout $S$-morphisme
$f : Y \to X$ et tout voisinage infinitésimal $Y'$ de $Y$ sur $S$ auquel $f$
se prolonge (localement) en $f' : Y' \to X$, la stratification fournit un
système \emph{transitif} d'isomorphismes entre les $f'^{*}\underline{F}$, pour
les divers prolongements $f'$. La donnée d'une stratification revient donc à
celle, pour tout tel $Y'$, d'un \emph{prolongement} de $\underline{F}$ en un
Module sur $Y'$, avec des conditions de fonctorialité évidentes. La condition
de prolongement est automatique si $X$ est formellement lisse sur $S$.

\begin{tikzcd}
X \arrow[dd] & Y \arrow[l, "f"'] \arrow[d, hook] \\
 & Y' \arrow[ul, "f'"] \\
S &
\end{tikzcd}

\emph{4.2. Les sites.} D'où le site et le topos infinitésimaux de $X$ sur
$S$, et le site et le topos stratifiants, au sens des
conventions.\footnote{La page 54 donne les titres de ces définitions, non les
définitions ; voir la note des conventions. Le croquis de la page 53 a un
trait vertical sans flèche entre $X$ et $S$ ; nous mettons la flèche
structurale.} Ils coïncident si $X$ est formellement lisse sur $S$.
Attention : le site n'a pas d'objet qui couvre l'objet final, en
général.\footnote{« le site n'a \emph{pas de générateur d'objet final} ! »,
page 54. C'est ce que contourneront, aux pages 59 et 62, le faisceau
$\widetilde{X}$ pour le site stratifiant et le choix d'un plongement dans un
$Y$ formellement lisse pour le site infinitésimal.} On décrit les faisceaux
sur ces sites par des faisceaux ordinaires, de même les faisceaux d'anneaux
et de Modules. Les faisceaux \emph{spéciaux} sont ceux dont les homomorphismes
de changement de base sont des isomorphismes — aujourd'hui, les
\emph{cristaux} ; et les Modules spéciaux sur le site stratifiant de $X/S$
sont la même chose que les Modules stratifiés sur $X/S$.

\emph{4.3. Un théorème et une conjecture}, énoncés pour être détaillés plus
tard.

\emph{Théorème.} Si $S$ est de caractéristique nulle et $X$ lisse sur $S$,
\[
  H^{*}\bigl((X/S)_{\mathrm{inf}}, \mathcal{O}\bigr) \simeq
  H^{*}\bigl((X/S)_{\mathrm{strat}}, \mathcal{O}\bigr) \simeq
  H^{*}_{\mathrm{DR}}(X/S) .
\]
\footnote{Sous la flèche, deux lignes encadrées et biffées, de lecture très
incertaine, portent sur $Rf_{*}$. Le théorème est le lemme de Poincaré
formel des pages 13 et 14, une fois la cohomologie stratifiante identifiée à
celle du site.}

\emph{Conjecture.} Si $X$ est localement de type fini sur $\mathbb{C}$,
\[
  H^{*}\bigl((X/\mathbb{C})_{\mathrm{strat}}, \mathcal{O}\bigr)
  \simeq H^{*}(X^{\mathrm{an}}, \mathbb{C}) ,
\]
et de même pour le site infinitésimal. Pour $X$ lisse, elle résulte du
théorème et du théorème de comparaison de la page 21 : on obtient ainsi, en
caractéristique nulle, une description de la cohomologie de De Rham
\emph{sans utiliser de différentielles}. Elle est intéressante pour $X$
singulier, où les formes différentielles échouent.\footnote{Pour le site
infinitésimal, l'énoncé est aujourd'hui un théorème : par un plongement
local dans un schéma lisse $Y$, la cohomologie infinitésimale se calcule par
le complété formel de $\Omega^{\cdot}_Y$ le long de $X$, et Hartshorne
(« On the De Rham cohomology of algebraic varieties », \emph{Publ.\ Math.\
IHÉS} 45, 1975) montre que celui-ci calcule $H^{*}(X^{\mathrm{an}},
\mathbb{C})$. Pour le site stratifiant d'un $X$ singulier, nous ne savons pas
dire où en est la question ; elle rejoint le « Pb » des pages 70 et 71.}

\pagerange{56}{58}
\section*{La reprise du 6 décembre 1966 (pages 56 à 58)}

Soit $\mathcal{F}$ un Module quasi-cohérent stratifié sur $X/S$. Il définit un
Module sur $(X/S)_{\mathrm{strat}}$, et un faisceau cosimplicial sur
$X_{\mathrm{zar}}$, son complexe de Čech--Alexander
\[
  C^{\nu}(\mathcal{F}) = \varprojlim_{i} \mathcal{F}_{\Delta^{(\nu)}(i)} ,
\]
défini par l'image inverse le long de l'une quelconque des $\nu+1$
projections $\Delta^{(\nu)}(i) \to X$, toutes canoniquement isomorphes grâce
à la stratification.\footnote{La page 56 dit « n'importe laquelle des $n$
projections » pour le produit à $n+1$ facteurs ; il y en a $n+1$.} On a
\[
  \boxed{\;H^{*}\bigl((X/S)_{\mathrm{strat}}, \mathcal{F}\bigr) \simeq
  \mathbb{H}^{*}\bigl(X_{\mathrm{zar}}, C^{*}(\mathcal{F})\bigr)\;}
\]
On a donc deux théories cohomologiques : pour $\mathcal{F}$ quasi-cohérent,
les $H^{i}(X_{\mathrm{zar}}, \mathcal{F})$ ; pour $\mathcal{F}$ stratifié,
les $H^{i}((X/S)_{\mathrm{strat}}, \mathcal{F})$. La seconde s'exprime par la
première, comme hypercohomologie zariskienne d'un complexe dont les
différentielles ne sont \emph{pas linéaires}. L'inverse ne semble pas vrai
\emph{a priori} ; il l'est pourtant, à peu de chose près.

\emph{La linéarisation.} À tout Module quasi-cohérent $\mathcal{F}$, sans
stratification, on associe un pro-Module pro-quasi-cohérent stratifié
$\mathcal{Q}(\mathcal{F})$, dont la cohomologie stratifiante est la
cohomologie zariskienne \emph{ordinaire} de $\mathcal{F}$, et qui est
fonctoriel non seulement en les homomorphismes linéaires, mais en les
opérateurs différentiels relatifs à $S$. Dans le cas très particulier où la
diagonale $X \to X \times_S X$ est une immersion nilpotente, stratification
et donnée de descente relativement à $S$ se confondent, et
$\mathcal{Q}(\mathcal{F}) = f^{*}f_{*}\mathcal{F}$. L'intérêt : \emph{l'hypercohomologie
zariskienne ordinaire d'un complexe d'opérateurs différentiels s'interprète
comme la cohomologie stratifiante d'un complexe de faisceaux stratifiés, à
différentielles linéaires.} C'est le programme que réalisent, l'été suivant,
les pages 4 à 9.

En termes de parties principales : un opérateur différentiel $M \to N$
induit un morphisme $\mathcal{P}^{\nu}[M] \to \mathcal{P}^{\nu}[N]$, qui n'est
pas $\mathcal{P}^{\nu}$-linéaire, mais seulement $\mathcal{P}^{\nu-1}$-linéaire
($\nu \geq 1$) quand on voit $\mathcal{P}^{\nu}[M]$ comme
$\mathcal{P}^{\nu-1}$-Module par restriction des scalaires le long de
$\alpha^{\nu-1}$ ; vu ainsi, on le note $\mathcal{Q}^{\nu-1}[M]$. En
particulier $\mathcal{Q}^{0}[M]$ est un foncteur de $\mathrm{Mod}(X/S)$ vers
les pro-$\mathcal{O}_X$-Modules, et $\alpha$ fait de $\mathcal{Q}^{*}$ une
algèbre cosimpliciale sur $\mathcal{P}^{*}$, avec
\[
  \boxed{\;\mathcal{Q}^{\nu-1}[M] = \mathcal{P}^{\nu} \otimes_{\mathcal{O}_X} M\;}
\]
\footnote{La formule encadrée du bas de la page 58 a un indice lu « $i-1$ »
ou « $i$ » ; seule la lecture $\mathcal{Q}^{i-1}[M] = \mathcal{P}^{i} \otimes
M$ s'accorde avec la phrase qui la précède et avec la page 5. La longue note
verticale de la marge de la page 58 est illisible, sauf le mot
« Anti-Riesz », que nous n'interprétons pas.}

\pagerange{59}{60}
\section*{Le plan, et la cohomologie stratifiante par Čech (pages 59 à 61)}

Le plan de la page 59 ordonne toute la série : §\,1 cohomologie de De Rham
ordinaire ; §\,2 tour d'horizon $\ell$-adique et motivations ; §\,3
Washnitzer--Monsky ; §\,4 stratifications, et la relation avec Gauss--Manin ;
§\,5 site cristallin et site stratifiant, un théorème et une conjecture en
caractéristique nulle ; §\,6 calculs « čechistes », en caractéristique
quelconque ; §\,7 comparaison de la cohomologie de De Rham et de la
cohomologie cristalline ; §\,8 retour sur la caractéristique $p$.\footnote{La
numérotation des feuillets antérieurs ne suit pas exactement ce plan : la
section Gauss--Manin y porte le numéro §\,3 (page 46), celle des topos
§\,4 (page 53). Le §\,6 est le seul que la suite écrive sous ce numéro ; voir
le fil du dossier pour les §\,7 et §\,8.}

\emph{6.1. Cohomologie stratifiante.} Le topos stratifiant de $X/S$ a un
objet $\widetilde{X}$, représenté par $X$ lui-même, qui \emph{couvre l'objet
final} — c'est ici que sert la restriction aux objets qui se rétractent sur
$X$. D'où une suite spectrale de Leray
\[
  E_{2}^{pq} = H^{p}\bigl(\nu \mapsto H^{q}(\widetilde{X}^{\nu+1}, F)\bigr)
  \Longrightarrow H^{p+q}\bigl((X/S)_{\mathrm{strat}}, F\bigr) ,
\]
avec $\widetilde{X}^{\nu+1} = \varinjlim_{i} \widetilde{\Delta^{(\nu)}(i)}$, et
pour tout objet $(U, X')$ du site, $H^{*}\bigl((U, X'), F\bigr) \simeq
H^{*}(X', F_{(U, X')})$. Supposons $X$ affine, donc les $\Delta^{(\nu)}(i)$
affines, et les $F_{\Delta^{(\nu)}(i)}$ quasi-cohérents, à flèches de
transition surjectives en $i$. Alors $H^{q}(\widetilde{\Delta^{(\nu)}(i)}, F) =
0$ pour $q > 0$, le système des $H^{0}$ est strict, et par Mittag-Leffler
\[
  H^{q}(\widetilde{X}^{\nu+1}, F) = 0 \ (q > 0) ,
  \qquad
  H^{0}(\widetilde{X}^{\nu+1}, F) = \varprojlim_{i} F\bigl(\Delta^{(\nu)}(i)\bigr) ,
\]
d'où l'isomorphisme canonique
\[
  \boxed{\;H^{*}\bigl((X/S)_{\mathrm{strat}}, F\bigr) \simeq
  H^{*}\bigl(\nu \mapsto F(X^{\nu+1}/X)\bigr)\;}
  \qquad (X \text{ affine}) ,
\]
par exemple pour $F = \mathcal{O}$ : la cohomologie stratifiante d'un affine
est la cohomologie du complexe des fonctions sur les complétés formels des
diagonales.\footnote{Le second membre de l'hypothèse est lu « et
indépendants mutuellement », sous réserve, aux pages 60 et 61. La condition
de Mittag-Leffler, que la page nomme (« la M.L. ») sans la formuler, est
assurée par la surjectivité des flèches de transition, que nous
écrivons : elle a lieu pour $F = \mathcal{O}$, les $\Delta^{(\nu)}(i) \to
\Delta^{(\nu)}(i+1)$ étant des immersions fermées entre affines.}

\pagerange{61}{61}
\subsection*{Le cas non affine (page 61)}

Si $X$ n'est plus supposé affine, les
$U \mapsto F(U^{\nu+1}/U) = \varprojlim_{i} F\bigl(\Delta^{(\nu)}_{U}(i)\bigr)$
sont des faisceaux $\mathcal{F}^{(\nu)}$ sur $X_{\mathrm{zar}}$, qui forment un
faisceau cosimplicial $\mathcal{F}^{*}$ — le complexe de Čech--Alexander des
conventions —, et
\[
  H^{*}\bigl((X/S)_{\mathrm{strat}}, F\bigr) \simeq
  \mathbb{H}^{*}(X_{\mathrm{zar}}, \mathcal{F}^{*}) ,
  \qquad
  \boxed{\;E_{2}^{pq} = H^{p}\bigl(\nu \mapsto H^{q}(X_{\mathrm{zar}},
  \mathcal{F}^{(\nu)})\bigr) \Longrightarrow
  H^{p+q}\bigl((X/S)_{\mathrm{strat}}, F\bigr)\;}
\]
ce qui généralise la formule de la page 56.

\pagerange{61}{64}
\section*{Le site infinitésimal par un plongement (pages 61 à 64)}

\emph{6.2.} Pour le site infinitésimal, $\widetilde{X}$ ne couvre plus l'objet
final quand $X$ n'est pas formellement lisse sur $S$. Supposons qu'on ait une
immersion
\[
  X \hookrightarrow Y , \qquad Y \text{ formellement lisse sur } S ;
\]
c'est toujours possible localement sur $X$, et globalement si $X$ est
quasi-projectif sur $S$.\footnote{La page 62 est lacunaire ici ;
« localement sur $X$ » est biffé puis l'idée reprise à la page 64, où le
passage aux ouverts affines est fait explicitement. Le croquis de la page 62
— $X \hookrightarrow Y$ au-dessus de $S$ — est remplacé par cette ligne.}
Soit $Y(i)$ le $i$-ème voisinage infinitésimal de $X$ dans $Y$. Le faisceau
« représenté » par $(X, Y)$ est $\widetilde{Y} = \varinjlim_i
\widetilde{Y(i)}$, et la lissité formelle de $Y$ assure qu'il couvre l'objet
final : tout épaississement d'un ouvert de $X$ se prolonge localement à
$Y$. D'où
\[
  E_{2}^{pq} = H^{p}\bigl(\nu \mapsto H^{q}(\widetilde{Y}^{\nu+1}, F)\bigr)
  \Longrightarrow H^{p+q}\bigl((X/S)_{\mathrm{inf}}, F\bigr) .
\]
Supposons $F$ quasi-cohérent sur chaque objet, à flèches de transition
surjectives pour les immersions $X'' \hookrightarrow X'$ — c'est le cas d'un
cristal en Modules quasi-cohérents —, et $X$ affine, donc les $Y(i)$
affines. Comme $H^{*}(\widetilde{X'}, F) = H^{*}(X'_{\mathrm{zar}}, F_{X'})$,
on trouve encore $E_2^{pq} = 0$ pour $q \neq 0$, et
\[
  H^{*}\bigl((X/S)_{\mathrm{inf}}, F\bigr) \simeq
  H^{*}\bigl(\nu \mapsto F(Y^{\nu+1}/X)\bigr) ,
\]
où $Y^{\nu+1}/X$ est le complété formel de $Y^{\nu+1}$ le long de $X$
plongé diagonalement.\footnote{L'ajout de la page 63, « et simultanément
bien si $X'' \hookrightarrow X'$ ! », est ce que nous rendons par la
surjectivité des flèches de transition ; la lecture est incertaine. La page
dit aussi « supposons $U$ dans $X$ affines », que nous lisons : $X$ affine.}
Si $X$ n'est pas affine, on introduit les faisceaux sur $X_{\mathrm{zar}}$
\[
  \mathcal{F}^{(\nu)}(U) = F\bigl(V^{\nu+1}/U\bigr) ,
\]
$V$ un ouvert de $Y$ qui induit $U$ ; ils forment un faisceau cosimplicial
$\mathcal{F}^{*}$, et
\[
  H^{*}\bigl((X/S)_{\mathrm{inf}}, F\bigr) \simeq
  \mathbb{H}^{*}(X_{\mathrm{zar}}, \mathcal{F}^{*}) ,
  \qquad
  E_{2}^{pq} = H^{p}\bigl(\nu \mapsto H^{q}(X, \mathcal{F}^{(\nu)})\bigr)
  \Longrightarrow H^{p+q}\bigl((X/S)_{\mathrm{inf}}, F\bigr) .
\]
C'est le complexe de Čech--Alexander d'un plongement, sous le nom qu'il porte
encore aujourd'hui.

\pagerange{64}{65}
\section*{Invariance par les immersions nilpotentes (pages 64 et 65)}

\emph{6.3.} Soit $X_0 \hookrightarrow X$ une immersion \emph{nilpotente}, et
$F$ sur $(X/S)_{\mathrm{inf}}$ provenant par restriction d'un $F_0$ sur
$(X_0/S)_{\mathrm{inf}}$. Pour tout $\nu$, les complétés formels
$Y^{\nu+1}/X_0$ et $Y^{\nu+1}/X$ sont les mêmes, puisqu'ils ne dépendent que
du fermé sous-jacent ; donc $\mathcal{F}^{(\nu)} \simeq
\mathcal{F}_0^{(\nu)}$ sur $X_{\mathrm{zar}} = X_{0,\mathrm{zar}}$, et
\[
  H^{*}\bigl((X/S)_{\mathrm{inf}}, F\bigr) \xrightarrow{\ \sim\ }
  H^{*}\bigl((X_0/S)_{\mathrm{inf}}, F_0\bigr) .
\]
Établi pour un plongement global dans un $Y$ formellement lisse, cela s'étend
au cas général par un recouvrement de $X$ par des ouverts affines. Il
resterait à voir si l'on peut se passer de l'hypothèse de quasi-cohérence
faite sur $F$.\footnote{C'est l'\emph{invariance topologique} du site
infinitésimal. Elle est propre au site sans puissances divisées : pour le
site cristallin à puissances divisées, elle ne vaut que pour les immersions
compatibles aux puissances divisées.}

Supposons en particulier $X$ \emph{lisse} sur $S$, et $S_0 \hookrightarrow
S$ une immersion nilpotente ; alors $X_0 = X \times_S S_0 \to X$ en est une,
et, en regardant $X_0$ comme schéma \emph{sur $S$},
\[
  H^{*}\bigl((X/S)_{\mathrm{inf}}, \mathcal{O}\bigr) \simeq
  H^{*}\bigl((X_0/S)_{\mathrm{inf}}, \mathcal{O}\bigr) .
\]
\emph{Le premier membre ne dépend donc pas, à isomorphisme canonique près, de
la façon dont on a relevé la situation lisse $X_0/S_0$ en une situation lisse
sur $S$.} Si $S$ est de caractéristique nulle, il est aussi canoniquement
isomorphe à $H^{*}(X_{\mathrm{zar}}, \Omega^{\cdot}_{X/S})$ (théorème de la
page 54) ; on a donc une définition de la cohomologie de De Rham pour un
morphisme quelconque $f : X \to S$ de base de caractéristique nulle, et
mieux, on devrait pouvoir définir $Rf_{*}(\Omega^{\cdot}_{X/S})$ comme un
\emph{cristal absolu} sur $S$, c'est-à-dire l'étendre à tout voisinage
infinitésimal de $S$, et pas seulement à ceux qui se rétractent sur
$S$.\footnote{C'est ici l'analogue exact, en caractéristique nulle, de
l'invariance que Washnitzer et Monsky affirmaient en caractéristique $p$ :
l'indépendance du relèvement est un théorème pour le site infinitésimal en
caractéristique nulle. Que la page 52 vienne de constater que l'analogue naïf
échoue en caractéristique $p$ est ce qui fait de ce passage le pivot du
dossier.}

\pagerange{66}{68}
\section*{L'anneau cosimplicial de Čech--Alexander (pages 66 à 68)}

\emph{6.4.} Pour $F = \mathcal{O}$, sur le site stratifiant ou infinitésimal,
$\mathcal{F}^{*}$ est un \emph{anneau cosimplicial} sur $X_{\mathrm{zar}}$.
Dans le cas du plongement de 6.2, c'est la limite projective des anneaux
cosimpliciaux $\mathcal{F}(i)^{*}$ relatifs aux $Y(i)$ (qui ont tous l'espace
de $X$) ; son étude se ramène donc essentiellement au contexte
stratifiant.

Chaque $\mathcal{F}^{(\nu)}$ est une $\mathcal{O}_X$-algèbre de $\nu+1$
façons, par les $\nu+1$ projections de $X^{\nu+1}/X$ sur $X$. Si l'on en
choisit une, $\mathrm{pr}_1$ par exemple, $\mathcal{F}^{(\nu)}$ devient
pro-quasi-cohérent, et les opérations cosimpliciales — donc les cobords —
doivent être vus comme des \emph{opérateurs différentiels d'ordre infini}.
C'est ainsi que les complexes d'opérateurs différentiels s'introduisent dans
la théorie de façon vraiment naturelle. Analogie avec les cochaînes de
Čech--Alexander.\footnote{La fin de 6.4, page 66, est en partie illisible ;
elle continue page 68, après une page 67 qui ne porte qu'un premier départ
de 6.5, biffé (« Soit $X$ un espace analytique »). L'analogie est celle des
cochaînes d'Alexander--Spanier, fonctions sur $X^{\nu+1}$ au voisinage de la
diagonale ; ici le voisinage est formel.}

\pagerange{68}{70}
\section*{Le cas complexe (pages 68 à 70)}

\emph{6.5.} Soit $S = \operatorname{Spec} \mathbb{C}$. La suite spectrale de
6.1 s'écrit
\[
  E_{2}^{pq} = H^{p}\bigl(\nu \mapsto H^{q}(X^{\nu+1}/X,
  \mathcal{O})\bigr) \Longrightarrow
  H^{p+q}\bigl((X/\mathbb{C})_{\mathrm{strat}}, \mathcal{O}\bigr) ,
\]
avec, au second membre, une « cohomologie formelle ». Si $X$ est propre, les
$H^{q}(\Delta^{(\nu)}(i), \mathcal{O})$ sont de dimension finie, et
Mittag-Leffler donne
\[
  H^{q}(X^{\nu+1}/X, \mathcal{O}) \simeq \varprojlim_{i}
  H^{q}\bigl(\Delta^{(\nu)}(i), \mathcal{O}\bigr) ,
  \qquad
  E_{2}^{pq} = \varprojlim_{i} H^{p}\Bigl(\nu \mapsto
  H^{q}\bigl(\Delta^{(\nu)}(i), \mathcal{O}\bigr)\Bigr) .
\]
Par GAGA, la cohomologie des schémas propres $\Delta^{(\nu)}(i)$ s'interprète
analytiquement, d'où
\[
  H^{*}\bigl((X/\mathbb{C})_{\mathrm{strat}}, \mathcal{O}\bigr) \simeq
  H^{*}\bigl(X^{\mathrm{an}}_{\mathrm{strat}}, \mathcal{O}\bigr)
  := \mathbb{H}^{*}\bigl(X^{\mathrm{an}}, \mathcal{F}^{*}_{X^{\mathrm{an}}}\bigr) ,
\]
où, pour un espace analytique $\mathfrak{X}$, $\mathcal{F}^{*}_{\mathfrak{X}}$
est le faisceau cosimplicial $U \mapsto \Gamma\bigl(U^{\nu+1}/U,
\mathcal{O}\bigr)$, $U$ parcourant les ouverts de $\mathfrak{X}$.

\emph{Conjecture.} Pour tout espace analytique complexe $\mathfrak{X}$, le
complexe $\mathcal{F}^{*}_{\mathfrak{X}}$ est, par l'augmentation évidente, une
résolution du faisceau constant $\mathbb{C}_{\mathfrak{X}}$. Cela semble
plausible, en regardant les « fibres formelles » des faisceaux.\footnote{Pour
$\mathfrak{X}$ lisse, c'est une forme du lemme de Poincaré formel. Pour
$\mathfrak{X}$ singulier, l'énoncé voisin où l'on complète un plongement
$\mathfrak{X} \subset \mathfrak{Y}$ lisse — le complété formel de
$\Omega^{\cdot}_{\mathfrak{Y}}$ le long de $\mathfrak{X}$ résout
$\mathbb{C}_{\mathfrak{X}}$ — est démontré par Bloom et Herrera (« De Rham
cohomology of an analytic space », \emph{Invent.\ Math.} 7, 1969). Nous ne
savons pas si la conjecture telle qu'énoncée ici, sans plongement, a été
établie.}

\emph{Théorème.} Soit $X$ propre sur $\mathbb{C}$. Si la conjecture est vraie
pour $X^{\mathrm{an}}$, on a un isomorphisme canonique
\[
  H^{*}\bigl((X/\mathbb{C})_{\mathrm{strat}}, \mathcal{O}\bigr) \simeq
  H^{*}(X^{\mathrm{an}}, \mathbb{C}) .
\]
C'est immédiat : l'hypercohomologie d'une résolution de $\mathbb{C}$ est la
cohomologie de $\mathbb{C}$.

Supposons de plus $X$ contenu dans un schéma lisse $Y$. Alors l'homomorphisme
\[
  H^{*}\bigl((X/\mathbb{C})_{\mathrm{inf}}, \mathcal{O}\bigr) \longrightarrow
  H^{*}\bigl((X/\mathbb{C})_{\mathrm{strat}}, \mathcal{O}\bigr)
\]
est également un isomorphisme — sous la même hypothèse, puisque le second
membre est celui du théorème.\footnote{Le raisonnement de la page 70 est
elliptique. Il invoque, pour le premier membre, $H^{*}_{\mathrm{inf}}(X) =
\varprojlim_i H^{*}_{\mathrm{inf}}(Y(i))$ avec des flèches de transition
bijectives — ce qui résulte de 6.3, mais ne compare rien au site
stratifiant. Une hypothèse qui aurait fait marcher l'argument, la conjecture
pour les $Y(i)^{\mathrm{an}}$, est écrite puis biffée. Ce qui rend l'énoncé
vrai, sous la conjecture pour $X^{\mathrm{an}}$, est que le premier membre
est lui aussi isomorphe à $H^{*}(X^{\mathrm{an}}, \mathbb{C})$ — c'est le
théorème de Hartshorne cité à la page 55, par le complexe de 6.2 —, et que
la flèche est compatible aux deux identifications, ce qui reste à vérifier.
Nous énonçons donc le résultat, et non la démonstration de la page.} Cela
rejoint, dit-il, la conjecture de l'exposé précédent, celle de la page 55.

\pagerange{70}{71}
\section*{Un problème, et une remarque non archimédienne (pages 70 et 71)}

\emph{Problème.} Soit $X$ plat et localement de présentation finie sur $S$,
et $F$ sur $(X/S)_{\mathrm{inf}}$ satisfaisant aux hypothèses habituelles
(par exemple provenant d'un Module stratifié). A-t-on
\[
  H^{*}\bigl((X/S)_{\mathrm{inf}}, F\bigr) \xrightarrow{\ \sim\ }
  H^{*}\bigl((X/S)_{\mathrm{strat}}, F\bigr) \ ?
\]
C'est équivalent au problème suivant : pour $X$ comme ci-dessus et
$X_0 \to X$ une immersion nilpotente, a-t-on
$H^{*}((X/S)_{\mathrm{strat}}, F) \simeq H^{*}((X_0/S)_{\mathrm{strat}}, F)$ ?
— l'invariance topologique de 6.3, mais pour le site stratifiant. On peut
supposer $S$, $X$, $X_0$ affines.\footnote{« plat » est lu sous réserve. Le
dossier pose le problème et n'y revient pas. L'équivalence se voit ainsi :
le site infinitésimal est invariant par immersion nilpotente (6.3) et
coïncide avec le site stratifiant pour $X$ lisse ; plonger localement $X$
dans un lisse et comparer aux voisinages infinitésimaux ramène l'un à
l'autre.}

\emph{Remarques.} Si $S = \operatorname{Spec} k$, $k$ un corps valué complet
non archimédien, il y a lieu de se poser la question de la conjecture de la
page 69 pour les espaces analytiques « flasques » sur $k$ ; mais rien
n'indique globalement qu'elle soit vraisemblable. Le dossier s'arrête
là.\footnote{Un passage barré et hachuré mentionnait d'abord les espaces
rigides-analytiques de Tate. La fin de la page est en partie illisible ;
« flasques » est entre guillemets sur la page, et nous ne savons pas quelle
classe d'espaces il désigne.}

\end{document}
